\documentclass[12pt,a4paper]{article}

\usepackage[margin=1in]{geometry}
\usepackage{booktabs}
\usepackage{longtable}
\usepackage{listings}
\usepackage{xcolor}
\usepackage{hyperref}
\usepackage{titlesec}
\usepackage{parskip}
\usepackage{fancyhdr}
\usepackage{graphicx}
\usepackage{array}
\usepackage{tabularx}
\usepackage{amsmath}

% Colors
\definecolor{codebg}{rgb}{0.95,0.95,0.95}
\definecolor{keywordcolor}{rgb}{0.13,0.29,0.53}
\definecolor{commentcolor}{rgb}{0.25,0.5,0.35}
\definecolor{stringcolor}{rgb}{0.73,0.13,0.13}

% Code listing style
\lstset{
  backgroundcolor=\color{codebg},
  basicstyle=\ttfamily\small,
  keywordstyle=\color{keywordcolor}\bfseries,
  commentstyle=\color{commentcolor}\itshape,
  stringstyle=\color{stringcolor},
  breaklines=true,
  breakatwhitespace=true,
  frame=single,
  rulecolor=\color{gray!50},
  framesep=4pt,
  xleftmargin=8pt,
  xrightmargin=8pt,
  showstringspaces=false,
  tabsize=4,
  captionpos=b,
}

% Section formatting
\titleformat{\section}{\large\bfseries}{}{0em}{}[\titlerule]
\titleformat{\subsection}{\normalsize\bfseries}{}{0em}{}

% Header / footer
\pagestyle{fancy}
\fancyhf{}
\lhead{\small DAS 839 -- NoSQL Systems}
\rhead{\small Phase 1 \& 2 Progress Report}
\cfoot{\thepage}

\hypersetup{
  colorlinks=true,
  linkcolor=keywordcolor,
  urlcolor=keywordcolor,
}

% -----------------------------------------------------------------------------
\begin{document}

\begin{titlepage}
  \centering
  \vspace*{2cm}
  {\LARGE\bfseries DAS 839 -- NoSQL Systems\par}
  \vspace{0.5cm}
  {\large End-Semester Project --- Phase 1 \& 2 Progress Report\par}
  \vspace{2cm}
  {\large\bfseries Ashok\par}
  \vspace{0.5cm}
  {\normalsize April 2026\par}
  \vspace{3cm}
  \hrule
  \vspace{0.4cm}
  {\small
    Dataset: NASA HTTP Server Log, July 1995 \quad|\quad
    1,891,714 records, 196 MB
  }
  \vspace{0.4cm}
  \hrule
\end{titlepage}

\tableofcontents
\newpage

% -----------------------------------------------------------------------------
\section{Overview}

Both the \textbf{MapReduce} and \textbf{MongoDB} pipelines are fully implemented and tested.
Each pipeline reads the same raw NASA HTTP server log file, applies the same parsing rules,
runs the same three analytical queries (Q1--Q3), loads results into shared PostgreSQL tables,
and produces identical output --- verified side by side on the same 1,000-record test slice.

The system is designed around a single CLI entry point (\texttt{main.py}) that accepts a
\texttt{--pipeline} flag, allowing any pipeline to be swapped in without touching any other
component.

% -----------------------------------------------------------------------------
\section{Project Structure}

\begin{lstlisting}[language=bash, caption={Repository layout}]
final_proj/
+-- main.py                          # CLI entry point (mapreduce + mongodb)
+-- config.py                        # Central constants (Hadoop, MongoDB, PostgreSQL)
+-- setup.sql                        # PostgreSQL schema (run once)
+-- requirements.txt                 # Python dependencies
+-- data/
|   \-- NASA_access_log_Jul95        # 1.89M records, 196 MB
+-- pipelines/
|   +-- mapreduce/
|   |   +-- parser.py                # Shared log parser (also used by MongoDB)
|   |   +-- runner.py                # MapReduce orchestrator
|   |   +-- loader.py                # HDFS output -> PostgreSQL
|   |   \-- mr_jobs/
|   |       +-- q1_mapper.py / q1_reducer.py
|   |       +-- q2_mapper.py / q2_reducer.py
|   |       \-- q3_mapper.py / q3_reducer.py
|   \-- mongodb/
|       +-- runner.py                # MongoDB orchestrator
|       \-- loader.py                # Aggregation pipelines -> PostgreSQL
+-- reporting/
|   \-- reporter.py                  # Reads PostgreSQL, prints formatted report
\-- db/
    \-- connection.py                # psycopg2 connection factory
\end{lstlisting}

% -----------------------------------------------------------------------------
\section{Infrastructure}

\begin{table}[h!]
\centering
\begin{tabularx}{\textwidth}{>{\bfseries}lX}
\toprule
Component & Details \\
\midrule
Hadoop    & 3.3.6 at \texttt{/home/ashok\_ubun/hadoop/}, single-node pseudo-distributed \\
HDFS      & \texttt{hdfs://localhost:9000}, replication = 1 \\
YARN      & Running with \texttt{mapreduce\_shuffle} auxiliary service \\
MongoDB   & 8.0, running on \texttt{localhost:27017}, data at \texttt{/var/lib/mongodb/} \\
PostgreSQL & 16.13, custom cluster at \texttt{/home/ashok\_ubun/pgdata\_nasa/}, port 5433 \\
Python    & 3.12 with \texttt{psycopg2-binary} and \texttt{pymongo} \\
Dataset   & \texttt{NASA\_access\_log\_Jul95} --- 1,891,714 lines, 196 MB \\
\bottomrule
\end{tabularx}
\caption{Infrastructure components}
\end{table}

% -----------------------------------------------------------------------------
\section{Data Model and Parsing}

\subsection{Log Format}

Raw log lines follow the NCSA Combined Log Format:

\begin{lstlisting}
199.72.81.55 - - [01/Jul/1995:00:00:01 -0400] "GET /history/apollo/ HTTP/1.0" 200 6245
\end{lstlisting}

\subsection{Parsed Fields}

\texttt{parser.py} uses a single regular expression to extract 8 structured fields per line:

\begin{table}[h!]
\centering
\begin{tabularx}{\textwidth}{lX}
\toprule
Field & Description \\
\midrule
\texttt{host}              & Client IP address or hostname \\
\texttt{log\_date}         & Date of the request (\texttt{DD/Mon/YYYY}) \\
\texttt{log\_hour}         & Hour of the request (0--23) \\
\texttt{http\_method}      & HTTP verb (\texttt{GET}, \texttt{POST}, etc.) \\
\texttt{resource\_path}    & Requested URI path \\
\texttt{protocol\_version} & HTTP protocol version \\
\texttt{status\_code}      & HTTP response code \\
\texttt{bytes\_transferred}& Bytes sent; \texttt{-} mapped to \texttt{0} \\
\bottomrule
\end{tabularx}
\caption{Parsed log fields}
\end{table}

Unparseable lines return \texttt{None} and are counted as malformed --- they are never silently
dropped. The same \texttt{parse\_line()} function is imported by both the MapReduce and MongoDB
runners, ensuring zero parsing divergence between pipelines.

% -----------------------------------------------------------------------------
\section{PostgreSQL Schema}

All pipelines write to the same five tables, distinguished by \texttt{run\_id} and
\texttt{pipeline\_name}:

\begin{itemize}
  \item \textbf{\texttt{pipeline\_runs}} --- one row per run; stores pipeline name, start/end
        timestamps, total runtime, batch count, avg batch size, and malformed count.
  \item \textbf{\texttt{batch\_log}} --- one row per batch; records per-batch record and malformed
        counts.
  \item \textbf{\texttt{q1\_daily\_traffic}} --- Query 1 results: daily traffic aggregated by
        status code.
  \item \textbf{\texttt{q2\_top\_resources}} --- Query 2 results: top 20 most-requested resources
        with rank.
  \item \textbf{\texttt{q3\_hourly\_errors}} --- Query 3 results: hourly error analysis with error
        rate.
\end{itemize}

% -----------------------------------------------------------------------------
\section{Analytical Queries}

All three queries run on the same parsed data regardless of pipeline.

\subsection{Q1 -- Daily Traffic Summary}

Groups records by \texttt{(log\_date, status\_code)} and aggregates:
\begin{itemize}
  \item Total request count
  \item Total bytes transferred
\end{itemize}

\subsection{Q2 -- Top 20 Requested Resources}

Groups records by \texttt{resource\_path} and aggregates:
\begin{itemize}
  \item Total request count
  \item Total bytes transferred
  \item Count of distinct requesting hosts
\end{itemize}
Results are sorted by request count descending; the top 20 are ranked 1--20 in Python before
being written to PostgreSQL.

\subsection{Q3 -- Hourly Error Analysis}

Groups records by \texttt{(log\_date, log\_hour)} and aggregates:
\begin{itemize}
  \item Total request count
  \item Error count (HTTP status 400--599)
  \item Error rate (\%)
  \item Distinct hosts that generated at least one error
\end{itemize}

% -----------------------------------------------------------------------------
\section{Pipeline 1 -- MapReduce (Hadoop Streaming)}

\subsection{Architecture}

The MapReduce pipeline uses Hadoop Streaming to run Python mapper/reducer scripts on data
staged in HDFS.

\textbf{Flow:}
\begin{enumerate}
  \item \textbf{Preflight} --- confirms HDFS is reachable and the input file exists.
  \item \textbf{Create run record} --- inserts a row into \texttt{pipeline\_runs}, obtains
        \texttt{run\_id}.
  \item \textbf{Batch staging} --- reads the input file in chunks of \texttt{batch\_size} lines;
        parses each chunk; writes valid records as a TSV file to HDFS; logs each batch to
        \texttt{batch\_log}.
  \item \textbf{MapReduce jobs} --- runs Q1, Q2, Q3 Hadoop Streaming jobs sequentially on the
        staged HDFS data.
  \item \textbf{Load results} --- reads \texttt{part-*} output from HDFS via
        \texttt{hdfs dfs -cat} and inserts into PostgreSQL.
  \item \textbf{Finalize} --- updates \texttt{pipeline\_runs} with runtime, total records,
        malformed count, and average batch size.
\end{enumerate}

\subsection{Mapper / Reducer Scripts}

\begin{table}[h!]
\centering
\begin{tabularx}{\textwidth}{lX}
\toprule
Script & Behaviour \\
\midrule
\texttt{q1\_mapper.py}  & Emits \texttt{log\_date \textbackslash t status\_code \textbackslash t 1 \textbackslash t bytes} \\
\texttt{q1\_reducer.py} & Groups by \texttt{(date, status)}; sums request count and total bytes \\
\texttt{q2\_mapper.py}  & Emits \texttt{resource\_path \textbackslash t host \textbackslash t 1 \textbackslash t bytes} \\
\texttt{q2\_reducer.py} & Groups by path; sums count/bytes; collects distinct hosts in a set \\
\texttt{q3\_mapper.py}  & Emits \texttt{log\_date \textbackslash t log\_hour \textbackslash t host \textbackslash t is\_error \textbackslash t 1} \\
\texttt{q3\_reducer.py} & Groups by \texttt{(date, hour)}; counts errors/total; computes error\_rate; distinct error hosts \\
\bottomrule
\end{tabularx}
\caption{MapReduce scripts}
\end{table}

The key Hadoop Streaming flag \texttt{-D stream.num.map.output.key.fields=2} is used for Q1
and Q3 to enable grouping by a 2-field composite key.

\subsection{Run Command}

\begin{lstlisting}[language=bash]
export JAVA_HOME=/usr/lib/jvm/java-21-openjdk-amd64
python3 main.py --pipeline mapreduce \
    --batch-size 50000 \
    --input data/NASA_access_log_Jul95
\end{lstlisting}

% -----------------------------------------------------------------------------
\section{Pipeline 2 -- MongoDB Aggregation}

\subsection{Architecture}

The MongoDB pipeline inserts parsed documents directly into a MongoDB collection and runs
server-side aggregation pipelines to compute Q1--Q3.

\textbf{Flow:}
\begin{enumerate}
  \item \textbf{Preflight} --- pings MongoDB with \texttt{serverSelectionTimeoutMS=3000};
        checks input file exists.
  \item \textbf{Create run record} --- same \texttt{INSERT} into \texttt{pipeline\_runs} as
        MapReduce.
  \item \textbf{Create indexes} --- three compound indexes on \texttt{log\_records} before
        insert to accelerate aggregations.
  \item \textbf{Batch ingest} --- reads in \texttt{batch\_size} chunks; parses with
        \texttt{parse\_line()}; stamps each document with \texttt{run\_id}; calls
        \texttt{insert\_many(ordered=False)} per batch; logs each batch to \texttt{batch\_log}.
  \item \textbf{Aggregations} --- calls \texttt{loader.load\_q1/q2/q3()}.
  \item \textbf{Finalize} --- same \texttt{UPDATE} to \texttt{pipeline\_runs} as MapReduce.
\end{enumerate}

\subsection{Aggregation Pipelines}

\begin{itemize}
  \item \textbf{Q1:} \texttt{\$match} (by run\_id) \textrightarrow{}
        \texttt{\$group (log\_date, status\_code)} \textrightarrow{} \texttt{\$sort}
  \item \textbf{Q2:} \texttt{\$match} \textrightarrow{} \texttt{\$group (resource\_path)}
        with \texttt{\$addToSet(host)} \textrightarrow{} \texttt{\$sort DESC} \textrightarrow{}
        \texttt{\$limit 20}; rank assigned in Python
  \item \textbf{Q3:} \texttt{\$match} \textrightarrow{} \texttt{\$group (log\_date, log\_hour)}
        with \texttt{\$cond} for error counting and \texttt{\$\$REMOVE} trick for distinct
        error hosts in a single pass
\end{itemize}

All aggregations use \texttt{allowDiskUse=True} (required for 1.89M documents).

\subsection{Run Command}

\begin{lstlisting}[language=bash]
python3 main.py --pipeline mongodb \
    --batch-size 50000 \
    --input data/NASA_access_log_Jul95
\end{lstlisting}

% -----------------------------------------------------------------------------
\section{Reporting}

\texttt{reporter.py} reads all results from PostgreSQL for a given \texttt{run\_id} regardless
of which pipeline produced them and prints a formatted report:

\begin{itemize}
  \item \textbf{Run metadata:} pipeline name, start/end timestamps, runtime, batch count,
        avg batch size, malformed record count.
  \item \textbf{Q1 table:} \texttt{log\_date | status\_code | request\_count | total\_bytes},
        sorted by date and status code.
  \item \textbf{Q2 table:} \texttt{rank | resource\_path | request\_count | total\_bytes |
        distinct\_hosts}, sorted by rank (1--20).
  \item \textbf{Q3 table:} \texttt{log\_date | log\_hour | error\_count | total\_count |
        error\_rate\% | distinct\_error\_hosts}, sorted by date and hour.
\end{itemize}

\begin{lstlisting}[language=bash]
python3 main.py --report --run-id <run_id>
\end{lstlisting}

% -----------------------------------------------------------------------------
\section{Test Run Results}

Both pipelines were tested on the same 1,000-record slice (\texttt{batch\_size=200}).

\begin{table}[h!]
\centering
\begin{tabularx}{\textwidth}{lcccccccc}
\toprule
Pipeline   & Run ID & Records & Malformed & Batches & Runtime & Q1 rows & Q2 rows & Q3 rows \\
\midrule
MapReduce  & 3 & 1,000 & 0 & 5 & $\sim$71s & 4 & 20 & 1 \\
MongoDB    & 5 & 1,000 & 0 & 5 & 0.04s    & 4 & 20 & 1 \\
\bottomrule
\end{tabularx}
\caption{Test run comparison (1,000 records, batch\_size=200)}
\end{table}

Q1, Q2, and Q3 results are \textbf{identical} across both pipelines --- same rows, same values.

The 1,775$\times$ runtime difference between the pipelines reflects the fundamental architectural
difference: MongoDB processes data in-process on already-ingested documents, while MapReduce
incurs per-job JVM startup, HDFS I/O, and shuffle overhead even for small inputs. This gap
is expected to narrow (in proportional terms) at full scale.

% -----------------------------------------------------------------------------
\section{How to Run}

\subsection{Start Services}

\begin{lstlisting}[language=bash]
# Hadoop (HDFS + YARN)
export JAVA_HOME=/usr/lib/jvm/java-21-openjdk-amd64
/home/ashok_ubun/hadoop/sbin/start-dfs.sh
/home/ashok_ubun/hadoop/sbin/start-yarn.sh

# PostgreSQL (custom cluster, port 5433)
/usr/lib/postgresql/16/bin/pg_ctl -D /home/ashok_ubun/pgdata_nasa \
    -o "-p 5433 -k /home/ashok_ubun/pgdata_nasa/socket" \
    -l /home/ashok_ubun/pgdata_nasa/pg.log start

# MongoDB
sudo mongod --dbpath /var/lib/mongodb \
    --logpath /var/log/mongodb/mongod.log --fork
\end{lstlisting}

\subsection{Run the Schema Setup (once)}

\begin{lstlisting}[language=bash]
psql -p 5433 -d nasa -f setup.sql
\end{lstlisting}

\subsection{Run a Pipeline}

\begin{lstlisting}[language=bash]
cd /home/ashok_ubun/studies_ubun/nosql/final_proj

# MapReduce
export JAVA_HOME=/usr/lib/jvm/java-21-openjdk-amd64
python3 main.py --pipeline mapreduce \
    --batch-size 50000 --input data/NASA_access_log_Jul95

# MongoDB
python3 main.py --pipeline mongodb \
    --batch-size 50000 --input data/NASA_access_log_Jul95
\end{lstlisting}

\subsection{View Report}

\begin{lstlisting}[language=bash]
python3 main.py --report --run-id <run_id>
\end{lstlisting}

% -----------------------------------------------------------------------------
\section{What Comes Next (Phase 2)}

\begin{itemize}
  \item Add Apache Pig pipeline (\texttt{pipelines/pig/})
  \item Add Apache Hive pipeline (\texttt{pipelines/hive/})
  \item Run all 4 pipelines on both July and August datasets with the same batch size
  \item Compare: runtime, implementation complexity, batching behaviour, and reporting
        suitability
  \item Record a video demo showing live pipeline selection and report output
  \item Write a compact PDF report with architecture, design decisions, and comparative
        analysis across all four pipelines
\end{itemize}

\end{document}
